\section{Discussion and conclusion}\label{sec:discussion}

The central results distinguish three properties of a reconstruction
representation. Algebraic distinctness means that its coordinates permit
functions unavailable to a specified comparison span. A bounded correction
means that a fitted residual head cannot exceed a stated magnitude under
the feature-norm assumption. Predictive usefulness means that the complete
procedure achieves lower loss under a specified evaluation distribution.
None of these properties implies the next one without additional
assumptions or evidence.

The principal experiment supports a contribution from reconstruction
departures beyond the shared locality and principal-coordinate features.
It does not establish the prespecified practical advantage specific to
KAHM across the comparison family. The matched RBF reconstruction
achieves much of the same benefit, and both boosting and the selected
non-KAHM procedure have lower aggregate point loss. Failure to establish
a margin does not establish equality of two methods. Likewise, the
irregular-bank example demonstrates existence of distinct raw spans,
not distinctness for every bank fitted in the experiment.

The correction bound provides a precise magnitude guarantee for the
specified decoder. It does not establish robustness of the full
history-to-feature map: signed normalization and the affine fallback may
still be sensitive to the input. The baseline and targets are not
uniformly bounded, so the theorem gives no uniform forecast-error bound.
Its applicability to other bounded features also limits a claim of
methodological novelty based on the bound alone.

The empirical scope is finite. Twenty-four independent repetitions
quantify variability under four fixed simulation laws; they are not
24 independent application domains. The observation maps, finite burn-in,
history length, horizons, model grids, and resource settings are part of
the tested procedures. The simultaneous intervals use a paired-mean
$t$ approximation, and the 3\% improvement and 5\% noninferiority margins
express research criteria rather than application-specific requirements.
The supporting natural-data and adaptive studies impose further limits
on extrapolation of the simulated gains. The experiments are not a
comprehensive comparison with neural, pretrained, long-context
multivariate, or modern adaptive forecasting methods.

Computation must also be interpreted at the level measured. Fitting
costs include the candidates evaluated during selection. Prediction time,
serialized state size, owned array storage, and whole-process peak memory
measure different resources and cannot be substituted for one another.
The measurements characterize the supplied implementations on the
recorded environment; they do not establish minimum algorithmic costs
or satisfy an unspecified deployment budget.

KAHM reconstruction can therefore supply useful features to a finite
forecasting decoder, with a value that depends on geometry, feature
parameterization, regularization, and target structure. The contribution
of the present study is to characterize these dependencies and evaluate
them against matched alternatives. A general forecasting advantage
specific to KAHM remains unestablished by the reported evidence.

\section*{Reproducibility and research transparency}

The anonymous saved-forecast supplement contains predictions and
evaluation targets for all 192 task suites, along with fitting targets
that determine the loss-normalization scales. Its standalone program,
\path{verify_saved_forecasts.py}, reconstructs eight numerical analysis
tables and all six primary contrasts, verifies array identities, and
checks that each policy forecast equals its recorded selected procedure.
The supplement includes all ten procedures and all declared conditions,
regardless of predictive performance. The companion \emph{Supplementary
Studies} provides the complete methods and results for the supporting
investigations summarized in Section~\ref{sec:supporting-evidence}.

The principal study's locally recorded source and input inventory
contains 239 files. The saved evidence contains 1,920 prediction arrays
and 9,584,640 forecast values. The recorded selected-state replay and
separately implemented calculations provide computational consistency
checks. Reconstructing scores from saved forecasts, replaying selected
states, and refitting all model candidates are distinct operations.
The anonymous numeric supplement performs the first; it does not refit
models or reproduce feature-bound calculations and resource timings.
The complete fitting specification and numerical conventions are given
in Appendix~\ref{app:reconstruction-results}. The prospective record was
local, informed by the supporting studies, and not publicly preregistered.
